% =====================================================================
%  Shapeshifter: A Declarative Language for Mass Spectrometric
%  Experiment Specification
%
%  Language specification. Self-contained: the grammar, type system,
%  and operational semantics are given in full, and every claim about
%  the language is a claim about the artefact defined here.
% =====================================================================
\documentclass[11pt,a4paper]{article}

\usepackage[utf8]{inputenc}
\usepackage[T1]{fontenc}
\usepackage{lmodern}
\usepackage[a4paper,margin=2.4cm]{geometry}
\usepackage{amsmath,amssymb,amsthm,mathtools}
\usepackage{bm}
\usepackage{booktabs}
\usepackage{array}
\usepackage{longtable}
\usepackage{enumitem}
\usepackage{microtype}
\usepackage{xcolor}
\usepackage{listings}
\usepackage{graphicx}
\usepackage[round,sort&compress]{natbib}
\usepackage[colorlinks=true,linkcolor=blue!45!black,
            citecolor=blue!45!black,urlcolor=blue!45!black]{hyperref}
\usepackage{cleveref}

% ---- Theorem-like environments ----
\newtheoremstyle{ssthm}{\topsep}{\topsep}{\itshape}{}{\bfseries}{.}{.5em}{}
\theoremstyle{ssthm}
\newtheorem{theorem}{Theorem}[section]
\newtheorem{lemma}[theorem]{Lemma}
\newtheorem{proposition}[theorem]{Proposition}
\newtheorem{corollary}[theorem]{Corollary}
\theoremstyle{definition}
\newtheorem{definition}[theorem]{Definition}
\newtheorem{rulename}[theorem]{Rule}
\newtheorem{example}[theorem]{Example}
\theoremstyle{remark}
\newtheorem{remark}[theorem]{Remark}
\newtheorem{notation}[theorem]{Notation}

\crefname{rulename}{Rule}{Rules}
\Crefname{rulename}{Rule}{Rules}

% ---- Listing style for Shapeshifter source ----
\definecolor{sskw}{RGB}{30,64,175}
\definecolor{sscom}{RGB}{22,101,52}
\definecolor{ssstr}{RGB}{159,18,57}
\definecolor{ssns}{RGB}{180,83,9}
\definecolor{ssbg}{RGB}{250,250,250}
\lstdefinelanguage{shapeshifter}{
  morekeywords={import,objective,instrument,dataset,phase,validate,target_list},
  morekeywords=[2]{lavoisier,instrument,partition,cells,msms,db,observe,purpose,acquire,transform,analyse},
  sensitive=true,
  morecomment=[l]{//},
  morestring=[b]",
  keywordstyle=\color{sskw}\bfseries,
  keywordstyle=[2]\color{ssns},
  commentstyle=\color{sscom}\itshape,
  stringstyle=\color{ssstr},
  basicstyle=\ttfamily\small,
  backgroundcolor=\color{ssbg},
  frame=single,framesep=4pt,xleftmargin=6pt,xrightmargin=4pt,
  breaklines=true,showstringspaces=false,columns=fullflexible,tabsize=2,
}
\lstset{language=shapeshifter}

% ---- Shorthand ----
\newcommand{\R}{\mathbb{R}}
\newcommand{\N}{\mathbb{N}}
\newcommand{\bnf}{\;::=\;}
\newcommand{\alt}{\;\mid\;}
\newcommand{\shk}[1]{\texttt{#1}}
\newcommand{\evalto}{\Downarrow}
\newcommand{\Env}{\rho}
\newcommand{\Kinds}{\kappa}
\newcommand{\Ord}{\omega}
\newcommand{\Wsp}{W}
\newcommand{\abs}[1]{\lvert#1\rvert}
\newcommand{\Prog}{P}
\newcommand{\Blocks}{\mathcal{B}}
\newcommand{\Vals}{\mathcal{V}}
\newcommand{\Kset}{\mathcal{K}}
% Semantic brackets, defined locally so no extra package is required.
\providecommand{\llbracket}{\ensuremath{[\![}}
\providecommand{\rrbracket}{\ensuremath{]\!]}}

\title{\textbf{Shapeshifter: A Declarative Language for\\[3pt]
Mass Spectrometric Experiment Specification}\\[6pt]
\large Grammar, Type Discipline, Operational Semantics,\\
and a Two-Stage Compilation Model}

\author{%
Kundai Farai Sachikonye\\
\small Technical University of Munich\\
\small \texttt{kundai.sachikonye@wzw.tum.de}}

\date{\today}

\begin{document}
\maketitle

% =====================================================================
\begin{abstract}
\noindent
A mass spectrometric experiment is normally specified twice: once as a
set of instrument-vendor settings entered through a graphical interface,
and again as an informal prose description in a methods section. Neither
artefact is executable, neither is diffable, and the correspondence
between them is maintained by hand. We present \emph{Shapeshifter}
(source extension \shk{.ss}), a declarative domain-specific language in
which an experiment is a single executable document. This paper is the
language's specification.

We give: (i)~a complete lexical and context-free grammar in EBNF for the
six block forms (\shk{objective}, \shk{instrument}, \shk{dataset},
\shk{target\_list}, \shk{phase}, \shk{validate}) and the value language
they contain;
(ii)~a \emph{kind} discipline in which every value carried by the runtime
is tagged with one of a fixed set of result kinds, and in which kind is
assigned by the producing operation rather than inferred from data shape,
so that an empty result retains the kind of the operation that produced
it; (iii)~a small-step operational semantics over configurations
$\langle \Env, \Kinds, \Ord, L\rangle$ of environment, kind map,
declaration order and log, from which we prove three properties ---
\emph{declaration-order determinism} (\Cref{thm:determinism}),
\emph{phase isolation} (\Cref{thm:isolation}), and \emph{workspace
monotonicity} (\Cref{thm:monotone}), the last being the statement that no
executed statement removes a binding, so the workspace is an append-only
record of the run; (iv)~a two-stage compilation model separating a
\textsc{compile} stage, which parses and performs structural checks
without executing any phase, from an \textsc{execute} stage, which runs
phases against a fixed AST --- and we prove that compile-stage
diagnostics are independent of runtime outcome
(\Cref{thm:stage-separation}), so that a program can be checked without
being run; and (v)~a standard library organised into nine namespaces, each operation
specified by its argument signature, result kind, and totality
behaviour. The library divides into a \emph{generative} half, which
synthesises records from declared parameters, and an \emph{acquisition}
half, which reads vendor-neutral \shk{mzML} and derives coordinates,
linkages, and coherence statistics from measured spectra; we show that
the two halves meet in a common record kind, so that a measured and a
predicted experiment are the same type of object and may be compared
without conversion (\Cref{thm:record-parity}).

We report a reference implementation written from this document alone
and a program executed against a public library of $5{,}328$ tandem
spectra. The run confirms stage separation in practice --- structural
checking costs $0.95$\,ms and opens no file, against $3.4$\,s and
$6.5$\,MB for execution --- and returns a verdict against the
coordinate transformation the program tests, failing three of four
criteria stated in the source before the run. We report that negative
outcome, and three under-specifications that writing the implementation
exposed, because a specification language whose worked examples all
succeed has not been shown to be capable of anything in particular.

Two design decisions are argued for explicitly and their consequences
proved. First, \emph{declaration order is the execution order}: phases
run in source order and statements within a phase run in source order,
with no dependency inference, which makes the document its own schedule.
Second, \emph{failure is local and named}: an unresolvable analyte class
is reported as a warning and skipped rather than aborting the run,
whereas an empty class set is an error, and we characterise exactly which
conditions belong to which category (\Cref{prop:failure-classification}).
We close with a worked corpus of seven programs spanning untargeted
lipidomics, targeted acquisition with per-class ranges, proteomics,
timing-cell compilation, partition addressing, domain-constrained
observation, and a complete acquisition-and-analysis run over measured
mzML, and we state what the language does not provide.
\end{abstract}

\noindent\textbf{Keywords:} domain-specific language; mass spectrometry;
experiment specification; operational semantics; EBNF grammar; kind
system; two-stage compilation; declarative workflow; reproducibility.

\tableofcontents

% =====================================================================
\section{Introduction}
\label{sec:intro}
% =====================================================================

\subsection{The specification gap}

Consider what must be true for a mass spectrometric experiment to be
repeated by a second laboratory. The instrument must be configured
identically in every respect that matters: analyser type, polarity,
collision energy, acquisition window, isolation width, gradient
programme. The analyte scope must be identical: which compound classes,
over which ranges of chain length and unsaturation, with which adducts.
The post-acquisition filtering must be identical: which retention window,
which intensity threshold, which charge states retained.

In current practice none of this is recorded in an executable form.
Instrument settings live in a vendor method file, which is binary,
version-locked, and not diffable. Analyte scope lives in a spreadsheet or
in the analyst's working memory. Filtering lives in a sequence of
mouse interactions with a processing package
\citep{pluskal2010mzmine,smith2006xcms,tsugawa2015msdial}. The published
methods section is a prose summary written after the fact, and the
minimum-information reporting standards that exist
\citep{taylor2007miape,sumner2007msi} specify what must be
\emph{disclosed}, not a language in which it can be \emph{stated}.

The consequence is a specification gap: the artefact that determines the
experiment and the artefact that describes it are different objects, and
nothing enforces their agreement. This is a recognised problem in
computational science generally
\citep{peng2011reproducible,sandve2013ten}, and the response elsewhere
has been executable workflow description
\citep{koster2012snakemake,ditommaso2017nextflow,goecks2010galaxy,%
amstutz2016cwl}. Those systems describe pipelines over files. They do not
describe experiments over instruments, and they have no vocabulary for
analytes, adducts, or acquisition windows.

\subsection{What Shapeshifter is}

Shapeshifter is a declarative language in which the specification of a
mass spectrometric experiment is a single executable document. A program
declares an objective, one or more instrument configurations, optional
target lists, an ordered sequence of phases containing the operations to
perform, and optional validation blocks. Running the program produces
records, timing-cell registries, or partition address tables, together
with a log of what was done.

The language is deliberately small. It has no user-defined functions, no
control flow, no recursion, and no mutable state beyond single
assignment. This is a design choice with a specific justification, given
in \Cref{sec:design}: an experiment specification that admits arbitrary
computation is no longer a specification, because reading it no longer
tells you what will happen. Everything Shapeshifter can express is
visible in the source without simulating its execution.

\subsection{What this paper provides}

This paper is the specification of the language, and it is self-contained
in the sense that every claim it makes is a claim about the grammar,
type discipline, and semantics defined here. Specifically:

\begin{enumerate}[label=(\roman*),leftmargin=2.2em,itemsep=2pt]
\item \Cref{sec:lexical} gives the lexical structure, including the
      significant-indentation rule and the bracket-continuation rule that
      allows a single logical line to span several physical lines.
\item \Cref{sec:grammar} gives the context-free grammar in EBNF
      \citep{backus1960notation,knuth1964backus} for all five block forms
      and the value language.
\item \Cref{sec:kinds} defines the kind discipline and proves kind
      preservation (\Cref{thm:kind-preservation}).
\item \Cref{sec:semantics} gives the operational semantics and proves
      determinism, phase isolation, and workspace monotonicity.
\item \Cref{sec:stages} defines the two-stage compilation model and
      proves stage separation.
\item \Cref{sec:acquisition} defines the acquisition layer: the
      \shk{dataset} block, the three acquisition namespaces, and the
      record-parity property joining them to the generative half.
\item \Cref{sec:stdlib} specifies the generative standard library.
\item \Cref{sec:programs} works six complete programs.
\item \Cref{sec:execution} reports a reference implementation and an
      executed experiment, including its negative outcome.
\item \Cref{sec:limits} states what the language does not do.
\end{enumerate}

\subsection{Relation to existing tools}

Shapeshifter is not a replacement for a data-processing package. It does
not read raw files, does not perform peak picking, and does not compute
false-discovery rates. Its relation to tools such as MZmine, XCMS,
MS-DIAL, OpenMS, or MaxQuant
\citep{pluskal2010mzmine,smith2006xcms,tsugawa2015msdial,%
rost2016openms,cox2008maxquant} is that of a specification language to an
implementation: those tools process the data an experiment produces;
Shapeshifter states what experiment is to be performed. The nearest
analogue in current practice is the targeted-method document of Skyline
\citep{maclean2010skyline}, which is likewise a declarative artefact
describing an acquisition; Shapeshifter differs in being a textual
language with a published grammar rather than an application document
format.

% =====================================================================
\section{Design rationale}
\label{sec:design}
% =====================================================================

Before the formal development we state the three decisions that shape the
language, because each is a restriction and each restriction is
deliberate.

\subsection{Declaration order is execution order}

Workflow languages commonly infer an execution order from a dependency
graph: a task runs when its inputs are available
\citep{koster2012snakemake,ditommaso2017nextflow}. Shapeshifter does not.
Phases execute in the order they appear in the source, and statements
within a phase execute in the order they appear within the phase.

The justification is that an experiment specification is read by people
who need to know the order of operations, and dependency inference makes
the order a derived property that must be computed to be known. A
document whose execution order is its reading order can be checked by
reading. The cost is that the author must order the phases correctly;
the language will not reorder them, and a statement referring to a
variable not yet bound will find it unbound rather than triggering
an out-of-order evaluation.

\subsection{No user-defined abstraction}

The language has no procedures, no macros, and no imports of user code.
The only callable operations are the fixed standard library of
\Cref{sec:stdlib}. This bounds what a program can do to what the
specification enumerates, which is what makes \Cref{prop:enumerable}
below --- that the set of possible effects of any program is finite and
computable from its source --- true at all. A language with
user-defined abstraction cannot have that property, and for an
experiment specification the property is worth more than the
expressiveness it costs \citep{mernik2005when,vandeursen2000domain}.

\subsection{Failure is local where the run can continue}

A specification that aborts on the first unrecognised token is
frustrating to use; a specification that silently ignores everything it
does not understand is dangerous. Shapeshifter takes an intermediate
position that is stated precisely in
\Cref{prop:failure-classification}: conditions that leave the run
meaningful are warnings and are named in the log; conditions that leave
the run meaningless are errors and halt it. An unknown analyte class is
the former --- the remaining classes still constitute an experiment. An
empty class set after filtering is the latter --- there is no experiment
left to run.

% =====================================================================
\section{Lexical structure}
\label{sec:lexical}
% =====================================================================

\subsection{Source text and lines}

\begin{definition}[Source, physical line]
\label{def:source}
A Shapeshifter \emph{source} is a finite sequence of Unicode characters.
A \emph{physical line} is a maximal substring containing no newline. Each
physical line $\ell$ carries an \emph{indent} $\mathrm{ind}(\ell)\in\N$,
the number of leading whitespace characters, and a \emph{body}
$\mathrm{body}(\ell)$, obtained by deleting any suffix beginning with the
comment marker \shk{//} and then trimming surrounding whitespace.
\end{definition}

\begin{definition}[Blank and comment lines]
\label{def:blank}
A physical line is \emph{vacuous} if $\mathrm{body}(\ell)$ is empty. This
covers both wholly blank lines and lines containing only a comment.
Vacuous lines are discarded before parsing and carry no syntactic
significance; in particular a vacuous line does not terminate a block.
\end{definition}

\begin{remark}[Why comments are stripped before indentation is used]
\label{rem:comment-first}
Comment stripping precedes vacuity testing, so a line consisting of
whitespace followed by \shk{//} is vacuous regardless of its indent. Were
the order reversed, an indented comment inside a block would be a
non-vacuous line at block indentation and would be offered to the field
parser, which would reject it.
\end{remark}

\subsection{Logical lines and bracket continuation}

Block structure is determined by indentation, but a single logical
statement frequently exceeds one physical line --- an array of object
literals describing a target panel, for example. Shapeshifter resolves
this with a bracket-balance rule rather than an explicit continuation
character.

\begin{definition}[Bracket balance]
\label{def:balance}
For a string $s$ let $\mathrm{op}(s)$ count occurrences of
$\shk{[}$, $\shk{(}$, $\shk{\{}$ and $\mathrm{cl}(s)$ count occurrences of
$\shk{]}$, $\shk{)}$, $\shk{\}}$. The \emph{balance} is
$\mathrm{bal}(s)=\mathrm{op}(s)-\mathrm{cl}(s)$.
\end{definition}

\begin{definition}[Logical line]
\label{def:logical}
Let $\ell_1,\ell_2,\dots$ be the non-vacuous physical lines in order. A
\emph{logical line} beginning at $\ell_i$ is formed as follows. If
$\mathrm{bal}(\mathrm{body}(\ell_i))\le 0$, the logical line is
$\ell_i$ alone. Otherwise successive bodies
$\mathrm{body}(\ell_{i+1}),\mathrm{body}(\ell_{i+2}),\dots$ are appended,
separated by a single space, accumulating balance, until the running
balance returns to zero or the source is exhausted. The logical line
inherits the indent of $\ell_i$.
\end{definition}

\begin{proposition}[Continuation terminates and is unambiguous]
\label{prop:continuation}
For any finite source, the construction in \Cref{def:logical} terminates,
and every non-vacuous physical line belongs to exactly one logical line.
\end{proposition}

\begin{proof}
Termination: each step consumes one physical line from a finite
sequence, so the inner loop runs at most as many times as there are
remaining lines. Partition: the construction proceeds left to right and
each step consumes a contiguous, non-empty run of lines beginning at the
current position; the next logical line begins immediately after. Hence
the logical lines partition the non-vacuous physical lines.
\end{proof}

\begin{remark}[Balance is counted on bodies, not raw lines]
\label{rem:balance-body}
Because balance is computed after comment stripping, a bracket appearing
inside a comment does not affect continuation. A bracket appearing inside
a string literal does, which is a known limitation recorded in
\Cref{sec:limits}.
\end{remark}

\subsection{Tokens}

\begin{definition}[Token classes]
\label{def:tokens}
The lexical classes are:
\begin{itemize}[leftmargin=1.6em,itemsep=1pt]
\item \textsc{Keyword}: one of \shk{import}, \shk{objective},
      \shk{instrument}, \shk{dataset}, \shk{phase}, \shk{validate},
      \shk{target\_list}.
\item \textsc{Ident}: a maximal run matching \shk{[A-Za-z\_][A-Za-z0-9\_]*}.
\item \textsc{QualName}: a dotted sequence of identifiers,
      \shk{[A-Za-z\_][A-Za-z0-9\_.]*}, used for namespaced operations.
\item \textsc{Number}: any string accepted by the host numeric reader,
      including decimal and scientific forms such as \shk{1.0e12}.
\item \textsc{String}: a run delimited by matching \shk{"} or \shk{'}.
\item \textsc{Bool}: \shk{true} or \shk{false}.
\item \textsc{Punct}: one of \shk{:}, \shk{,}, \shk{=},
      \shk{[}, \shk{]}, \shk{(}, \shk{)}, \shk{\{}, \shk{\}}.
\end{itemize}
\end{definition}

% =====================================================================
\section{Grammar}
\label{sec:grammar}
% =====================================================================

\subsection{Programs and blocks}

\begin{definition}[Grammar of programs]
\label{def:grammar}
A program is a sequence of imports and blocks:
\[
\begin{array}{rcl}
\textit{program}    &\bnf& \{\,\textit{import}\,\}\;\{\,\textit{block}\,\}\\[2pt]
\textit{import}     &\bnf& \shk{import}\ \textsc{QualName}\\[2pt]
\textit{block}      &\bnf& \textit{objective} \alt \textit{instrument}
                           \alt \textit{dataset} \alt \textit{targetlist}\\
                    &   & \alt\ \textit{phase} \alt \textit{validate}\\[4pt]
\textit{objective}  &\bnf& \shk{objective}\ \textsc{Ident}\ \shk{:}\ \textit{fieldblock}\\
\textit{instrument} &\bnf& \shk{instrument}\ \textsc{Ident}\ \shk{:}\ \textit{fieldblock}\\
\textit{dataset}    &\bnf& \shk{dataset}\ \textsc{Ident}\ \shk{:}\ \textit{fieldblock}\\
\textit{targetlist} &\bnf& \shk{target\_list}\ \textsc{Ident}\ \shk{:}\ \textit{fieldblock}\\
\textit{phase}      &\bnf& \shk{phase}\ \textsc{Ident}\ \shk{:}\ \textit{stmtblock}\\
\textit{validate}   &\bnf& \shk{validate}\ \textsc{Ident}\ \shk{:}\ \textit{stmtblock}
\end{array}
\]
\end{definition}

The six block forms divide into two groups by their body: \emph{field
blocks} carry declarative key--value data, and \emph{statement blocks}
carry executable statements. This division is the primary structural
distinction in the language.

\begin{definition}[Field and statement blocks]
\label{def:bodies}
Let $d$ be the indent of the block header. Then
\[
\begin{array}{rcl}
\textit{fieldblock} &\bnf& \{\,\textit{field}\,\}^{+}_{>d}\\
\textit{field}      &\bnf& \textsc{Ident}\ \shk{:}\ \textit{value}\\[4pt]
\textit{stmtblock}  &\bnf& \{\,\textit{stmt}\,\}^{+}_{>d}\\
\textit{stmt}       &\bnf& \textit{assign} \alt \textit{callstmt}\\
\textit{assign}     &\bnf& \textsc{Ident}\ \shk{=}\ \textit{expr}\\
\textit{callstmt}   &\bnf& \textsc{QualName}\ \shk{(}\ \textit{args}\ \shk{)}
\end{array}
\]
where $\{\,X\,\}^{+}_{>d}$ denotes one or more occurrences of $X$ on
logical lines of indent strictly greater than $d$.
\end{definition}

\begin{definition}[Expressions and arguments]
\label{def:expr}
\[
\begin{array}{rcl}
\textit{expr} &\bnf& \textsc{QualName}\ \shk{(}\ \textit{args}\ \shk{)}
                     \quad\text{(call)}\\
              &\alt& \textit{value} \quad\text{(literal or variable reference)}\\[4pt]
\textit{args} &\bnf& [\ \textit{arg}\ \{\ \shk{,}\ \textit{arg}\ \}\ ]\\
\textit{arg}  &\bnf& \textsc{Ident}\ \shk{:}\ \textit{value}
\end{array}
\]
All arguments are named; there are no positional arguments.
\end{definition}

\begin{remark}[Why arguments are named]
\label{rem:named-args}
Positional arguments make a call site unreadable without consulting the
callee's signature, and reorderings are silent errors. In a
specification language the call site must be self-describing, so
Shapeshifter admits only named arguments. The cost is verbosity at every
call.
\end{remark}

\subsection{The value language}

\begin{definition}[Values]
\label{def:values}
\[
\begin{array}{rcl}
\textit{value} &\bnf& \textsc{String} \alt \textsc{Number} \alt \textsc{Bool}\\
               &\alt& \shk{[}\ [\ \textit{value}\ \{\ \shk{,}\ \textit{value}\ \}\ ]\ \shk{]}
                      \quad\text{(array)}\\
               &\alt& \shk{\{}\ [\ \textit{pair}\ \{\ \shk{,}\ \textit{pair}\ \}\ ]\ \shk{\}}
                      \quad\text{(object)}\\
               &\alt& \textsc{Ident} \quad\text{(bare word / variable reference)}\\[4pt]
\textit{pair}  &\bnf& \textsc{Ident}\ \shk{:}\ \textit{value}
\end{array}
\]
\end{definition}

\begin{definition}[Value parsing function]
\label{def:parseval}
The function $\mathcal{P}:\textsc{String}\to\Vals$ is defined on a
trimmed input $s$ by the first applicable clause:
\begin{enumerate}[label=(V\arabic*),leftmargin=2.6em,itemsep=1pt]
\item if $s$ is empty, $\mathcal{P}(s)=\mathbf{null}$;
\item if $s$ is delimited by matching quotes, $\mathcal{P}(s)$ is the
      interior, taken literally;
\item if $s=\shk{true}$ or $s=\shk{false}$, $\mathcal{P}(s)$ is the
      corresponding boolean;
\item if $s$ is bracketed by \shk{[}\dots\shk{]}, $\mathcal{P}(s)$ is an
      array: if the interior begins with \shk{\{} it is parsed as a
      sequence of object literals, otherwise as a comma-separated
      sequence of values, in both cases splitting only at bracket depth
      zero;
\item if $s$ is bracketed by \shk{\{}\dots\shk{\}}, $\mathcal{P}(s)$ is an
      object formed by splitting the interior at depth-zero commas and
      each part at its first \shk{:};
\item if $s$ parses as a number, $\mathcal{P}(s)$ is that number;
\item otherwise $\mathcal{P}(s)=s$, a bare word.
\end{enumerate}
\end{definition}

\begin{lemma}[Depth-zero splitting is well defined]
\label{lem:split}
Let $s$ be a string and let $\mathrm{split}_0(s)$ be the sequence of
maximal substrings of $s$ separated by commas occurring at bracket depth
zero, where depth is incremented by $\shk{[}$, $\shk{(}$, $\shk{\{}$ and
decremented by $\shk{]}$, $\shk{)}$, $\shk{\}}$, scanning left to right.
Then $\mathrm{split}_0$ is total, and if $s$ is bracket-balanced then
concatenating the parts with commas recovers $s$.
\end{lemma}

\begin{proof}
Totality is immediate: the scan visits each character once and the
partition is defined by the positions at which depth is zero and the
character is a comma. For the recovery claim, note the split removes
exactly the depth-zero commas and no other character; reinserting them at
the same positions restores $s$.
\end{proof}

\begin{remark}[Bare words and the ambiguity they carry]
\label{rem:bareword}
Clause (V7) means that an unquoted, non-numeric token parses as itself.
This is what allows a target list to name an instrument block
(\shk{instrument: OrbitrapXL}) without quoting. It also means a
misspelled keyword inside a value position becomes a string rather than
an error. \Cref{sec:limits} records this.
\end{remark}

% =====================================================================
\section{The kind discipline}
\label{sec:kinds}
% =====================================================================

Shapeshifter does not have a type system in the sense of a static
calculus of value shapes \citep{pierce2002types,cardelli1985understanding}.
It has a weaker but load-bearing discipline: every binding produced by
executing a statement is tagged with a \emph{kind}, and the kind is
determined by the operation that produced it.

\subsection{Kinds}

\begin{definition}[Kind set]
\label{def:kinds}
The kind set is
\[
\Kset=\left\{
\begin{array}{l}
\shk{records},\ \shk{cells},\ \shk{addresses},\ \shk{subharmonics},\
\shk{tensor},\ \shk{impossible},\\
\shk{scans},\ \shk{coords},\ \shk{linkage},\ \shk{coherence},\
\shk{features},\\
\shk{scalar},\ \shk{list},\ \shk{object}
\end{array}
\right\}.
\]
The first eleven are \emph{operation kinds}, assigned by a specific
standard library operation; of these, \shk{scans}, \shk{coords},
\shk{linkage}, \shk{coherence} and \shk{features} are produced only by
the acquisition namespaces of \Cref{sec:acquisition}. The last three are
\emph{structural kinds}, assigned to literal bindings by inspection.
\end{definition}

\begin{definition}[Kind assignment]
\label{def:classify}
The function $\mathrm{kind}(f, v)$ assigns a kind to a value $v$ produced
by operation $f$ (with $f=\bot$ for a literal):
\[
\mathrm{kind}(f,v)=
\begin{cases}
\shk{records} & f\in\{\shk{run\_experiment},\shk{run\_proteomics},
                     \shk{sebd\_search},\shk{db.*}\}\\
\shk{cells} & f=\shk{cells.compile}\\
\shk{addresses} & f=\shk{partition.compute\_addresses}\\
\shk{subharmonics} & f=\shk{msms.phase\_coherence}\\
\shk{tensor} & f=\shk{msms.virtual\_tensor}\\
\shk{impossible} & f=\shk{msms.impossible\_ions}\\
\shk{scans} & f\in\{\shk{acquire.read\_mzml},\shk{acquire.filter\_scans}\}\\
\shk{coords} & f=\shk{transform.sentropy}\\
\shk{linkage} & f=\shk{acquire.link\_dda}\\
\shk{coherence} & f=\shk{analyse.phase\_lock}\\
\shk{features} & f=\shk{transform.extract\_features}\\
\shk{list} & f=\bot,\ v \text{ an array}\\
\shk{object} & f=\bot,\ v \text{ an object}\\
\shk{scalar} & \text{otherwise.}
\end{cases}
\]
\end{definition}

The essential property is that kind depends on $f$ and not on the
observed shape of $v$.

\begin{theorem}[Kind is producer-determined]
\label{thm:kind-preservation}
Let $f$ be an operation kind-assigning to $k\in\Kset$. Then for every
value $v$ returned by $f$, including $v=[\,]$, the binding receives kind
$k$. Consequently a downstream consumer selecting on kind cannot
distinguish an empty result of $f$ from a non-empty one by kind alone,
and cannot misclassify an empty result as a value of some other kind.
\end{theorem}

\begin{proof}
Immediate from \Cref{def:classify}: the first six clauses are conditioned
on $f$ alone and are checked before the shape-conditioned clauses, which
apply only when $f=\bot$. Hence for $f\ne\bot$ in the listed set the
value never reaches a shape test.
\end{proof}

\begin{corollary}[Empty results are not silently retyped]
\label{cor:empty-kind}
An operation returning zero records yields a binding of kind
\shk{records}, not of kind \shk{list}. A stage selecting the primary result
by kind preference therefore treats it as an empty experiment rather than
as an unrelated array.
\end{corollary}

\begin{remark}[Why this matters more than it appears to]
\label{rem:kind-why}
Shape-based classification is the obvious implementation and it is
wrong in exactly one case, which is the case that matters: an operation
that legitimately produced nothing. Under shape inference an empty
record set is indistinguishable from an empty list of user data, and a
downstream selector will either skip it or mislabel it. Producer-based
kinding costs nothing and removes the failure mode.
\end{remark}

\subsection{What the discipline does not do}

\begin{remark}[No argument checking]
\label{rem:no-argcheck}
The kind discipline governs results, not arguments. An operation
receiving an argument of unexpected shape will fail at runtime, not at
compile time. There is no signature check in the \textsc{compile} stage
beyond block structure. This is a deliberate scope limit, not an
oversight; extending the discipline to arguments is discussed in
\Cref{sec:limits}.
\end{remark}

% =====================================================================
\section{Operational semantics}
\label{sec:semantics}
% =====================================================================

\subsection{Configurations}

\begin{definition}[Configuration]
\label{def:config}
A \emph{configuration} is a quadruple
$\langle \Env, \Kinds, \Ord, L\rangle$ where
\begin{itemize}[leftmargin=1.6em,itemsep=1pt]
\item $\Env:\textsc{Ident}\rightharpoonup\Vals$ is the \emph{environment},
      a partial map from identifiers to values;
\item $\Kinds:\textsc{Ident}\rightharpoonup\Kset$ is the \emph{kind map};
\item $\Ord\in\textsc{Ident}^{*}$ is the \emph{declaration order}, a
      duplicate-free sequence;
\item $L$ is the \emph{log}, a sequence of records
      $(\mathit{level},\mathit{message})$ with
      $\mathit{level}\in\{\shk{info},\shk{warn},\shk{error}\}$.
\end{itemize}
The initial configuration is
$\langle\emptyset,\emptyset,\varepsilon,\varepsilon\rangle$.
\end{definition}

\subsection{Evaluation of expressions}

\begin{rulename}[Literal and variable reference]
\label{rule:lit}
\[
\frac{}{\langle\Env,e\rangle\evalto \mathcal{P}(e)}\ \textsc{(Lit)}
\qquad
\frac{v\in\mathrm{dom}(\Env)}{\langle\Env,v\rangle\evalto \Env(v)}\ \textsc{(Var)}
\]
\textsc{Var} applies during argument resolution: a named argument whose
value is a bare word naming a bound variable resolves to that variable's
value. Otherwise the bare word is retained as a string.
\end{rulename}

\begin{rulename}[Call]
\label{rule:call}
\[
\frac{\begin{array}{c}
a'=\{\,(k,\textsc{Resolve}(\Env,v)) : (k,v)\in a\,\}\qquad
\llbracket f\rrbracket(a',\Env,A)=v'
\end{array}}
{\langle\Env,\Kinds,\Ord,L\rangle,\ f(a)\ \evalto\ v'}
\ \textsc{(Call)}
\]
where $A$ is the enclosing AST (supplying \shk{instrument} and
\shk{target\_list} blocks by name), $\llbracket f\rrbracket$ is the
denotation of $f$ from \Cref{sec:stdlib}, and
$\textsc{Resolve}(\Env,v)=\Env(v)$ if $v$ is a bare word in
$\mathrm{dom}(\Env)$ and $v$ otherwise.
\end{rulename}

\begin{remark}[Argument resolution is one level deep]
\label{rem:resolve-depth}
$\textsc{Resolve}$ inspects the argument value itself, not values nested
inside arrays or objects. An array literal containing bare words that
name bound variables is passed through unresolved. This is a
deliberate restriction: nested resolution would make the meaning of an
array literal depend on the ambient environment, so that the same
literal denotes different data in different phases.
\end{remark}

\subsection{Execution of statements and blocks}

\begin{rulename}[Assignment]
\label{rule:assign}
\[
\frac{\langle\Env,\Kinds,\Ord,L\rangle,\ e\ \evalto\ v \qquad
      k=\mathrm{kind}(\mathrm{fn}(e),v)}
{\langle\Env,\Kinds,\Ord,L\rangle \xrightarrow{\ x=e\ }
 \langle\Env[x\mapsto v],\ \Kinds[x\mapsto k],\ \Ord\cdot x,\ L\cdot\lambda\rangle}
\]
where $\mathrm{fn}(e)=f$ if $e$ is a call $f(a)$ and $\bot$ otherwise,
$\Ord\cdot x$ appends $x$ if absent and is the identity if present, and
$\lambda$ is the assignment log entry.
\end{rulename}

\begin{rulename}[Phase and program]
\label{rule:phase}
\[
\frac{C_0 \xrightarrow{s_1} C_1 \xrightarrow{s_2}\cdots\xrightarrow{s_n} C_n}
{C_0 \xrightarrow{\ \shk{phase}\ p:\ s_1\dots s_n\ } C_n}
\qquad
\frac{C_0 \xrightarrow{p_1} C_1 \cdots \xrightarrow{p_m} C_m}
{C_0 \xrightarrow{\ \Prog\ } C_m}
\]
Phases execute in source order; statements within a phase execute in
source order. Validation blocks execute before all phases.
\end{rulename}

\subsection{Properties}

\begin{theorem}[Declaration-order determinism]
\label{thm:determinism}
Let $\Prog$ be a program and suppose every standard library operation
invoked by $\Prog$ is deterministic in its arguments. Then execution of
$\Prog$ from the initial configuration yields a unique final
configuration, and the sequence $\Ord$ is exactly the sequence of
assignment targets in source order with repetitions removed after their
first occurrence.
\end{theorem}

\begin{proof}
By induction on the number of statements. The base case is the initial
configuration. For the step, \Cref{rule:assign} is the only rule that
alters $\Env$, $\Kinds$ or $\Ord$, and it is a function of the current
configuration and the statement: the expression evaluates to a unique
$v$ by determinism of $\llbracket\cdot\rrbracket$ and of
$\textsc{Resolve}$; $k$ is a function of $(\mathrm{fn}(e),v)$ by
\Cref{def:classify}; and the updates are functional. \Cref{rule:phase}
composes these in a fixed order given by the source. Hence the final
configuration is unique. The claim about $\Ord$ follows because
$\Ord\cdot x$ appends exactly when $x\notin\Ord$.
\end{proof}

\begin{theorem}[Phase isolation]
\label{thm:isolation}
Phases share $\Env$ but do not scope it: a binding created in phase $p_i$
is visible in every phase $p_j$ with $j>i$, and in no phase $p_j$ with
$j<i$. There is no phase-local binding.
\end{theorem}

\begin{proof}
\Cref{rule:phase} threads a single configuration through the phases in
order, and \Cref{rule:assign} extends $\Env$ globally. No rule restricts
or restores $\Env$ at a phase boundary. Visibility forward is therefore
immediate. Visibility backward fails because $\Env$ at the time $p_j$
executes, $j<i$, does not contain the binding: by
\Cref{thm:determinism} the configuration entering $p_j$ is determined by
statements preceding it in source order, which exclude those of $p_i$.
\end{proof}

\begin{remark}[Phases are documentation, not scope]
\label{rem:phase-doc}
\Cref{thm:isolation} says phases carry no scoping semantics. Their role
is organisational: they name stages of an experiment and they structure
the log. A reader may rely on a phase boundary to mean ``this is a
distinct stage of the experiment'' and may not rely on it to mean
``these bindings are local''.
\end{remark}

\begin{theorem}[Workspace monotonicity]
\label{thm:monotone}
For any program and any two configurations $C$, $C'$ with
$C\xrightarrow{s} C'$ under \Cref{rule:assign},
$\mathrm{dom}(\Env)\subseteq\mathrm{dom}(\Env')$ and $\Ord$ is a prefix
of $\Ord'$. Consequently the workspace --- the sequence of bindings in
declaration order --- grows monotonically, and no executed statement
removes a binding or reorders the workspace.
\end{theorem}

\begin{proof}
$\Env'=\Env[x\mapsto v]$ extends the domain by at most $\{x\}$ and
removes nothing. $\Ord'=\Ord\cdot x$ either equals $\Ord$ (if $x$ already
occurs) or appends one element; in both cases $\Ord$ is a prefix. No
other rule modifies $\Env$ or $\Ord$.
\end{proof}

\begin{corollary}[Rebinding preserves position]
\label{cor:rebind}
If $x$ is assigned twice, the second assignment updates $\Env(x)$ and
$\Kinds(x)$ but leaves the position of $x$ in $\Ord$ at its first
occurrence. The workspace therefore reflects the order in which names
were \emph{introduced}, not the order in which they were last written.
\end{corollary}

\begin{proposition}[Effects are enumerable from source]
\label{prop:enumerable}
For any program $\Prog$, the set of standard library operations that
$\Prog$ can invoke is computable from the source by inspection, without
executing $\Prog$, and equals the set of \textsc{QualName} heads
appearing in call position.
\end{proposition}

\begin{proof}
By \Cref{def:expr} a call site names its operation by a literal
\textsc{QualName}; there is no expression form producing a callable
value, no user-defined abstraction, and no dynamic dispatch. Hence the
call graph is the set of syntactic call heads, obtainable by a single
pass over the AST.
\end{proof}

\subsection{Failure classification}

\begin{proposition}[Failure classification]
\label{prop:failure-classification}
Execution distinguishes two failure regimes:
\begin{enumerate}[label=\textup{(\roman*)},leftmargin=2.2em,itemsep=2pt]
\item \textbf{Local failure (warning).} A named entity that cannot be
      resolved but whose omission leaves the operation meaningful is
      skipped, a \shk{warn} record naming it is appended to $L$, and
      execution continues. Unknown analyte class keys are the canonical
      instance.
\item \textbf{Global failure (error).} A condition under which the
      operation has no meaningful result raises an exception, which
      terminates execution of the program and is reported as an
      \shk{error} record. An empty admissible class set after (i) has been
      applied is the canonical instance.
\end{enumerate}
No condition is silently ignored: every local failure appends a record
naming the skipped entity.
\end{proposition}

\begin{proof}
By inspection of the operation denotations in \Cref{sec:stdlib}: each
resolution step over a named collection either appends a \shk{warn} and
continues, or, where the collection would be empty, raises. The
no-silence claim follows because the skip and the log append occur in the
same step.
\end{proof}

\begin{remark}[The asymmetry is the point]
\label{rem:asymmetry}
A specification language that halts on the first unknown class makes
exploratory editing painful; one that discards unknown classes silently
makes a typo indistinguishable from an intended omission.
\Cref{prop:failure-classification} takes the third option: continue, but
name what was dropped. The log is then a complete record of the
difference between what was written and what was run.
\end{remark}

% =====================================================================
\section{The two-stage compilation model}
\label{sec:stages}
% =====================================================================

\subsection{Stages}

\begin{definition}[Compile stage]
\label{def:compile-stage}
$\textsc{Compile}(\mathit{src})$ returns a tuple
$(\mathit{ok}, A, \mathit{ir}, T, D)$ where $A$ is the AST or
$\bot$, $\mathit{ir}$ is a serialisation of $A$, $T$ is a sequence of
terminal lines each tagged \shk{stdout}, \shk{stderr} or \shk{stage}, and
$D$ is a sequence of diagnostics each tagged \shk{error} or
\shk{warning}. The stage performs: lexical analysis
(\Cref{sec:lexical}), parsing (\Cref{sec:grammar}), and the structural
checks of \Cref{def:structural}. \emph{It executes no phase.}
\end{definition}

\begin{definition}[Structural checks]
\label{def:structural}
The compile stage emits:
\begin{itemize}[leftmargin=1.6em,itemsep=1pt]
\item an \shk{error} diagnostic if parsing fails, with $\mathit{ok}=$ false
      and $A=\bot$;
\item a \shk{warning} if no \shk{objective} block is declared;
\item a \shk{warning} if no \shk{phase} block is declared, noting that
      nothing will execute.
\end{itemize}
\end{definition}

\begin{definition}[Execute stage]
\label{def:execute-stage}
$\textsc{Execute}(A)$ runs $A$ under \Cref{sec:semantics} and returns
$(\mathit{result}, T, L, \Wsp)$ where $\Wsp$ is the workspace,
$\mathit{result}$ is the primary result selected by
\Cref{def:primary}, and $T$ mirrors $L$ into the terminal stream.
\end{definition}

\begin{definition}[Primary result selection]
\label{def:primary}
The primary result is chosen by the first applicable clause:
\begin{enumerate}[label=(R\arabic*),leftmargin=2.6em,itemsep=1pt]
\item if $\Env(\shk{records})$ is a non-empty array, the result is that
      array with kind \shk{records};
\item if $\Env(\shk{registry})$ is an array, the result is that array with
      kind \shk{cells};
\item if $\Env(\shk{addresses})$ is an array, the result is that array with
      kind \shk{addresses};
\item otherwise, the first workspace entry whose kind lies in the
      preference order
      $\langle\shk{records},\shk{cells},\shk{addresses},\shk{subharmonics},
      \shk{tensor},\shk{impossible}\rangle$, or failing that the last
      workspace entry;
\item if the workspace is empty, the result has kind \shk{empty}.
\end{enumerate}
\end{definition}

\begin{remark}[Conventional names are a compatibility affordance]
\label{rem:conventional}
Clauses (R1)--(R3) privilege three conventional variable names. This is
a pragmatic concession to downstream consumers that expect them, not a
language-level scoping rule; (R4) is the general mechanism and is
sufficient on its own. An author who avoids the three names loses
nothing.
\end{remark}

\subsection{Stage separation}

\begin{theorem}[Stage separation]
\label{thm:stage-separation}
For any source $\mathit{src}$, the diagnostics $D$ returned by
$\textsc{Compile}(\mathit{src})$ are independent of
$\textsc{Execute}$: they are a function of $\mathit{src}$ alone.
Consequently a program may be checked for structural well-formedness
without invoking any standard library operation, and without any of the
effects those operations carry.
\end{theorem}

\begin{proof}
By \Cref{def:compile-stage} the stage's output depends on the token
sequence, the parse, and the three predicates of
\Cref{def:structural}, each of which is a property of $A$. No clause
consults $\Env$, $L$, or any operation denotation, and $\textsc{Compile}$
contains no application of \Cref{rule:call}. Hence $D$ factors through
$\mathit{src}\mapsto A$.
\end{proof}

\begin{corollary}[Checking is effect-free]
\label{cor:effect-free}
Structural checking of a program that would, on execution, query a
remote spectral database performs no query.
\end{corollary}

\begin{remark}[What stage separation does not give]
\label{rem:not-typecheck}
\Cref{thm:stage-separation} guarantees that checking is cheap and
effect-free. It does not guarantee that a program passing the compile
stage will execute without error: argument shapes are unchecked
(\Cref{rem:no-argcheck}), and unresolved names surface only at
\Cref{rule:call}. The compile stage certifies structure, not
executability.
\end{remark}

% =====================================================================
\section{Acquisition: from measured data to records}
\label{sec:acquisition}
% =====================================================================

The operations of \Cref{sec:stdlib} are \emph{generative}: they
synthesise records from declared parameters. A specification language
restricted to them can state what experiment to perform but cannot state
what to do with the data it produces, which leaves the larger half of
analytical practice outside the language. This section defines the
acquisition half: a \shk{dataset} block that names measured files, and
three namespaces that read them, transform them to coordinates, and
compute derived statistics.

The design constraint is that acquisition must not introduce a second
kind of program. A document that reads measured data and a document that
predicts an experiment must be the same sort of object, executed by the
same semantics, so that the two can appear in one file and their outputs
compared. \Cref{thm:record-parity} states the property that makes this
possible.

\subsection{The \shk{dataset} block}

\begin{definition}[Dataset block]
\label{def:dataset}
A \shk{dataset} block is a field block declaring a named collection of
input files together with the acquisition parameters under which they are
to be read:
\[
\begin{array}{rcl}
\textit{dataset} &\bnf& \shk{dataset}\ \textsc{Ident}\ \shk{:}\ \textit{fieldblock}
\end{array}
\]
The recognised fields are \shk{files} (an array of paths), \shk{format}
(default \shk{"mzML"}), \shk{vendor}, \shk{rt\_range}, \shk{polarity},
\shk{ms1\_threshold}, \shk{ms2\_threshold}, \shk{ms1\_precision},
\shk{ms2\_precision}, \shk{min\_peaks}, and \shk{dda\_top}. A
\shk{dataset} block declares data; it does not read it. Reading occurs
only when an \shk{acquire} operation names the block.
\end{definition}

\begin{remark}[Why declaration is separated from reading]
\label{rem:decl-vs-read}
Placing file paths and acquisition thresholds in a declarative block
rather than in call arguments has three consequences. The parameters are
stated once and shared by every operation that names the block, so a
threshold cannot silently differ between two phases. The declaration is
visible in the compile stage, so a program's inputs can be listed
without reading them (\Cref{prop:inputs-enumerable}). And the same
program can be repointed at a different cohort by editing one block.
\end{remark}

\begin{proposition}[Inputs are enumerable from source]
\label{prop:inputs-enumerable}
For any program $\Prog$, the set of input files $\Prog$ can read is
computable from the source without execution, and equals the union of
the \shk{files} fields of those \shk{dataset} blocks named by an
\shk{acquire} operation in call position.
\end{proposition}

\begin{proof}
By \Cref{def:dataset} file paths appear only in \shk{dataset} field
blocks, and by \Cref{def:acquire-ops} the only operations that open a
file take a \shk{dataset} name as an argument. Both the blocks and the
call sites are syntactic, so the set is obtained by one pass over the
AST, as in \Cref{prop:enumerable}.
\end{proof}

\begin{corollary}[Compile-stage input audit]
\label{cor:audit}
By \Cref{thm:stage-separation} and \Cref{prop:inputs-enumerable}, the
compile stage can report exactly which files a program will open, having
opened none of them.
\end{corollary}

\subsection{Namespace \shk{lavoisier.acquire}}

\begin{definition}[Acquisition operations]
\label{def:acquire-ops}
\hfill
\begin{longtable}{@{}p{0.26\textwidth}p{0.46\textwidth}p{0.20\textwidth}@{}}
\toprule
\textbf{Operation} & \textbf{Arguments} & \textbf{Kind / totality}\\
\midrule
\endhead
\shk{read\_mzml} &
\shk{dataset}, \shk{rt\_range?}, \shk{ms1\_threshold?}$=1000$,
\shk{ms2\_threshold?}$=10$, \shk{ms1\_precision?}, \shk{ms2\_precision?},
\shk{min\_peaks?}$=3$, \shk{vendor?}$=$\shk{"thermo"} &
\shk{scans}; partial\\
\addlinespace
\shk{filter\_scans} &
\shk{scans}, \shk{ms\_level?}, \shk{rt?}, \shk{mz?},
\shk{min\_peaks?}, \shk{polarity?} &
\shk{scans}; total\\
\addlinespace
\shk{link\_dda} &
\shk{scans}, \shk{tolerance?}$=0.01$, \shk{dda\_top?} &
\shk{linkage}; total\\
\addlinespace
\shk{xic} &
\shk{scans}, \shk{mz}, \shk{tolerance?} &
\shk{scans}; total\\
\bottomrule
\end{longtable}
\end{definition}

\begin{definition}[Scan record]
\label{def:scan}
A value of kind \shk{scans} is an array of \emph{scan records}, each
carrying: a scan identifier, an MS level, a retention time, a polarity,
an optional precursor $m/z$, and a peak list of $(m/z,\,\text{intensity})$
pairs. Thresholding and minimum-peak filtering are applied during
reading, so a scan record admitted to the array satisfies the declared
parameters by construction.
\end{definition}

\begin{remark}[Vendor neutrality is inherited, not claimed]
\label{rem:vendor}
\shk{read\_mzml} reads the open \shk{mzML} standard
\citep{martens2011mzml,deutsch2008mzml}. The \shk{vendor} field selects
among known controlled-vocabulary dialects and does not read
vendor-native formats. A program naming a vendor states how to interpret
the mzML it is given; it does not request a conversion.
\end{remark}

\begin{definition}[DDA linkage]
\label{def:linkage}
A value of kind \shk{linkage} is an array of \emph{DDA events}, each
pairing one MS1 survey scan with the MS2 scans acquired from it. An MS2
scan is assigned to the most recent preceding MS1 scan, and to a
precursor within \shk{tolerance} of its recorded precursor $m/z$. Scans
that cannot be assigned are retained in an \emph{unlinked} partition of
the result rather than discarded.
\end{definition}

\begin{remark}[Unlinked scans are reported, not dropped]
\label{rem:unlinked}
Retaining unlinked scans is an instance of
\Cref{prop:failure-classification}(i): the linkage is still meaningful
without them, so the run continues, but the count is named in the log and
the scans remain addressable. A pipeline that silently discarded them
would report a linkage rate of one by construction.
\end{remark}

\subsection{Namespace \shk{lavoisier.transform}}

\begin{definition}[Coordinate transformation operations]
\label{def:transform-ops}
\hfill
\begin{longtable}{@{}p{0.26\textwidth}p{0.46\textwidth}p{0.20\textwidth}@{}}
\toprule
\textbf{Operation} & \textbf{Arguments} & \textbf{Kind / totality}\\
\midrule
\endhead
\shk{sentropy} &
\shk{scans}, \shk{alpha?}, \shk{beta?}, \shk{k\_neighbors?},
\shk{precursor\_mz?} &
\shk{coords}; total\\
\addlinespace
\shk{extract\_features} &
\shk{coords} &
\shk{features}; total\\
\addlinespace
\shk{to\_records} &
\shk{coords}, \shk{linkage?} &
\shk{records}; total\\
\bottomrule
\end{longtable}
\end{definition}

The \shk{sentropy} operation assigns to each peak of each scan a triple
$(S_k, S_t, S_e)$. The components are computed as follows, with $m$ the
peak $m/z$ vector, $\hat I$ the intensity vector normalised to unit
maximum, $m^{\star}$ the precursor $m/z$ where known and $\max m$
otherwise, and $\alpha,\beta$ the declared weights:
\begin{align}
S_k &= -\log_2(\hat I + \epsilon) \;+\; \alpha\,\frac{m}{m^{\star}},
      \label{eq:sk}\\[2pt]
S_t &= \exp\!\left(-\beta\,\frac{\lvert m-\bar m\rvert}{\sigma_m}\right)
       \cdot g(t_R),
      \label{eq:st}\\[2pt]
S_e^{(i)} &= -\!\!\sum_{j\in N_k(i)}\! p_j\log_2 p_j,
       \qquad p_j=\frac{\hat I_j}{\sum_{l\in N_k(i)}\hat I_l},
      \label{eq:se}
\end{align}
where $\bar m$ and $\sigma_m$ are the mean and standard deviation of $m$,
$N_k(i)$ is the set of $k$ nearest peaks to peak $i$ in $m/z$, and
$g(t_R)$ is a retention-time modulation applied when a retention time is
available and the identity otherwise. The constant $\epsilon$ guards the
logarithm at zero intensity.

\begin{remark}[What the three components measure]
\label{rem:components}
\Cref{eq:sk} combines the Shannon self-information of a peak's normalised
intensity with a mass term, so $S_k$ is large for weak peaks and for
peaks near the precursor. \Cref{eq:st} is a Gaussian weighting about the
spectral centroid, so $S_t$ is large for peaks near the centre of the
$m/z$ distribution. \Cref{eq:se} is the Shannon entropy of the local
intensity distribution in a $k$-neighbourhood, so $S_e$ is large where
neighbouring peaks are of comparable intensity and small where one peak
dominates. All three are functions of the peak list alone: no library, no
formula assignment, and no instrument model enters.
\end{remark}

\begin{proposition}[Degenerate spectra are total, not exceptional]
\label{prop:degenerate}
\shk{sentropy} is total on scan records admitted under \Cref{def:scan}.
In particular: if $\sigma_m=0$ (all peaks at one $m/z$) then $S_t$ is the
constant $g(t_R)$; if a scan has a single peak then $N_k(i)=\{i\}$ and
$S_e=0$; and if no precursor is recorded then $m^{\star}=\max m$.
\end{proposition}

\begin{proof}
Each clause is a case of the definitions above, and each yields a defined
value rather than an error: the $\sigma_m=0$ branch replaces the
exponential by unity; the single-peak branch yields a one-element
neighbourhood whose normalised probability is $1$ and whose entropy is
therefore $0$; and the missing-precursor branch substitutes $\max m$,
which is defined because \Cref{def:scan} admits only scans carrying at
least one peak.
\end{proof}

\begin{remark}[Why totality here is worth stating]
\label{rem:totality-why}
Degenerate spectra are common in measured data --- a scan surviving
thresholding with a single peak above the cut is unremarkable. An
implementation that raised on them would fail partway through a cohort
for reasons unrelated to the analysis. \Cref{prop:degenerate} records
that each degenerate case has a defined value and says what it is, so the
behaviour is specified rather than incidental.
\end{remark}

\subsection{Namespace \shk{lavoisier.analyse}}

\begin{definition}[Analysis operations]
\label{def:analyse-ops}
\hfill
\begin{longtable}{@{}p{0.26\textwidth}p{0.46\textwidth}p{0.20\textwidth}@{}}
\toprule
\textbf{Operation} & \textbf{Arguments} & \textbf{Kind / totality}\\
\midrule
\endhead
\shk{phase\_lock} &
\shk{coords}, \shk{coherence\_threshold?}$=0.3$ &
\shk{coherence}; total\\
\addlinespace
\shk{summarise} &
\shk{coords} or \shk{records}, \shk{group\_by?} &
\shk{object}; total\\
\addlinespace
\shk{compare} &
\shk{a}, \shk{b}, \shk{metric?} &
\shk{object}; partial\\
\bottomrule
\end{longtable}
\end{definition}

\begin{definition}[Coherence]
\label{def:coherence}
\shk{phase\_lock} returns, per scan, a coherence statistic in $[0,1]$
together with the count of detections at or above
\shk{coherence\_threshold}. A value of kind \shk{coherence} carries both
the per-scan statistics and the threshold at which they were computed.
\end{definition}

\begin{remark}[The threshold travels with the result]
\label{rem:threshold-travels}
Carrying the threshold in the result rather than only in the call is
deliberate. A detection count is meaningless without the threshold that
produced it, and a report quoting one without the other cannot be
reproduced. Because the pair is a single value, a downstream consumer
cannot separate them by accident.
\end{remark}

\begin{remark}[A null detection is a result]
\label{rem:null-detection}
If no scan attains the threshold, \shk{phase\_lock} returns a
\shk{coherence} value with zero detections --- not an empty list and not
an error. By \Cref{thm:kind-preservation} the binding retains kind
\shk{coherence}, so the outcome is reported as ``no detections at
threshold $\theta$'' rather than as an absent result. Whether $\theta$
was appropriate is then a question the log makes askable.
\end{remark}

\subsection{The two halves meet}

\begin{theorem}[Record parity]
\label{thm:record-parity}
The generative operations of \Cref{sec:stdlib} and the acquisition
operation \shk{transform.to\_records} both yield values of kind
\shk{records}. Consequently a binding produced from measured data and a
binding produced from declared parameters are of the same kind, are
selected by the same clause of \Cref{def:primary}, and may be supplied to
the same downstream operation without conversion.
\end{theorem}

\begin{proof}
By \Cref{def:classify}, kind is a function of the producing operation;
both \shk{run\_experiment} and \shk{transform.to\_records} are listed as
assigning \shk{records}. \Cref{def:primary} selects on kind alone, and
operations consuming records are specified over the kind rather than over
the producer.
\end{proof}

\begin{corollary}[Measured and predicted experiments are comparable]
\label{cor:comparable}
A single program may bind one variable from a \shk{dataset} and another
from a virtual experiment and pass both to \shk{analyse.compare}, without
either being converted, because both are \shk{records}.
\end{corollary}

\begin{remark}[What parity does not assert]
\label{rem:parity-limits}
\Cref{thm:record-parity} is a statement about kinds, not about content.
It guarantees that the two halves compose; it does not guarantee that a
predicted record and a measured record are commensurable in any
scientific sense, and nothing in the language warrants that they are.
Whether a comparison is meaningful is a question for the analyst, and
\shk{analyse.compare} reports the metric it computed rather than a
verdict.
\end{remark}

\subsection{A complete acquisition program}

\begin{lstlisting}[caption={Reading a cohort, transforming, and analysing.}]
import lavoisier.acquire
import lavoisier.transform
import lavoisier.analyse

objective CohortSEntropy:
    target: "S-entropy coordinates for the M3/M4/M5 cohort"

dataset Cohort:
    files: [
        "A_M3_negPFP_03.mzML",
        "A_M4_posHILIC_01.mzML",
        "A_M5_negPFP_02.mzML"
    ]
    format: "mzML"
    rt_range: [0.5, 30.0]
    ms1_threshold: 1000
    ms2_threshold: 10
    min_peaks: 3

phase Read:
    scans = lavoisier.acquire.read_mzml(dataset: Cohort)
    ms1   = lavoisier.acquire.filter_scans(scans: scans, ms_level: 1)
    links = lavoisier.acquire.link_dda(scans: scans, tolerance: 0.01)

phase Transform:
    coords   = lavoisier.transform.sentropy(
        scans: ms1,
        alpha: 1.0,
        beta: 1.0,
        k_neighbors: 5
    )
    features = lavoisier.transform.extract_features(coords: coords)

phase Analyse:
    coherence = lavoisier.analyse.phase_lock(
        coords: coords,
        coherence_threshold: 0.3
    )
    summary   = lavoisier.analyse.summarise(coords: coords, group_by: "file")
\end{lstlisting}

Every parameter that determines the analysis --- the retention window,
both intensity thresholds, the minimum peak count, the transformation
weights $\alpha$ and $\beta$, the neighbourhood size $k$, and the
coherence threshold --- appears exactly once, in the source. By
\Cref{prop:inputs-enumerable} the three input files are listed by the
compile stage without being opened, and by \Cref{cor:audit} that listing
is effect-free.

\begin{example}[The threshold question, made askable]
\label{ex:threshold}
Suppose \shk{coherence} reports zero detections. By
\Cref{rem:null-detection} this is a reported outcome carrying the
threshold $0.3$, not an absent result. A reader of the document can see
that $0.3$ was chosen, can change it in one place, and can re-run; and a
reviewer can see that it was chosen rather than defaulted. Under the
alternative in which the threshold lives in a script and the report
quotes only the count, none of that is available.
\end{example}

\begin{remark}[Relation to fixed-stage pipelines]
\label{rem:stages}
The program above expresses in three phases what a procedural
implementation typically organises as a fixed sequence of stages ---
extraction, chromatography, ionisation, linkage, MS1 and MS2 analysis,
coordinate assignment, and validation. The correspondence is not exact
and is not meant to be: a fixed stage list is a property of one
implementation, whereas a phase sequence is written per experiment. What
the language fixes is not the stages but the discipline that each is
named, ordered, and parameterised in the source.
\end{remark}

% =====================================================================
\section{Standard library}
\label{sec:stdlib}
% =====================================================================

The generative library comprises twenty-two operations in seven
namespaces; the acquisition library of \Cref{sec:acquisition} adds ten
more in three further namespaces. We
specify each by argument signature, result kind, and totality.
Notation: $\shk{a?}$ marks an optional argument with its default;
\emph{total} means the operation returns for all well-formed arguments;
\emph{partial} means it may raise (\Cref{prop:failure-classification}(ii)).

\subsection{Namespace \shk{lavoisier.instrument}}

\begin{longtable}{@{}p{0.30\textwidth}p{0.44\textwidth}p{0.18\textwidth}@{}}
\toprule
\textbf{Operation} & \textbf{Arguments} & \textbf{Kind / totality}\\
\midrule
\endhead
\shk{run\_experiment} &
\shk{classes?}$=$\shk{["PC","PE"]}, \shk{polarity?}$=$\shk{"+"},
\shk{adducts?}, \shk{analyser?}$=$\shk{"orbitrap"},
\shk{collision\_energy?}$=25$, \shk{mz\_window?}$=[200,1500]$,
\shk{gradient\_min?}$=30$, \shk{filters?} &
\shk{records}; partial\\
\addlinespace
\shk{run\_proteomics} &
\shk{proteins?}$=$\shk{["HSA"]}, \shk{length\_min?}$=7$,
\shk{length\_max?}$=20$, \shk{mc\_max?}$=1$, \shk{polarity?},
\shk{analyser?}, \shk{collision\_energy?}$=28$, \shk{mz\_window?}$=[200,3000]$ &
\shk{records}; partial\\
\bottomrule
\end{longtable}

\begin{definition}[Class specification resolution]
\label{def:classspec}
The \shk{classes} argument admits two forms. A string entry \shk{"PC"}
denotes a class with default ranges. An object entry
$\{\shk{class}:c,\ \shk{carbons}:[x_0,x_1],\ \shk{db}:[y_0,y_1]\}$ denotes
explicit ranges. Resolution takes explicit ranges where given, otherwise
the class default, with the upper carbon bound capped at
$\min(X_{\max}^{\mathrm{def}}, X_{\min}^{\mathrm{def}}+8)$ to bound
enumeration. Unknown class keys are skipped with a warning; an empty
resolved set raises.
\end{definition}

\begin{definition}[Post-generation filters]
\label{def:filters}
The optional \shk{filters} object supports the keys \shk{mz}, \shk{rt},
\shk{db}, \shk{carbons} (each an inclusive interval $[lo,hi]$),
\shk{classes} and \shk{adducts} (each a membership list), and
\shk{min\_intensity} (a lower bound). A record is retained iff it
satisfies every present key. Absent keys impose no constraint.
\end{definition}

\begin{remark}[Filters are declarative and post hoc]
\label{rem:filters}
Filters are applied after generation, not folded into it. The record
count before and after filtering is logged, so the reduction is visible.
This is the property that makes a filtered experiment auditable: the
document states both the generated scope and the retained subset.
\end{remark}

\subsection{Namespace \shk{lavoisier.partition}}

\shk{compute\_addresses(lipids)} maps each input object
$\{\shk{name},\shk{mass},\shk{adduct}\}$ to a record carrying integer
coordinates $(n,\ell,m,s)$ derived from the mass, subject to
$0\le\ell\le n-1$ and $-\ell\le m\le\ell$. Kind \shk{addresses}; total.

\subsection{Namespace \shk{lavoisier.cells}}

\shk{compile(target\_list)} resolves the named \shk{target\_list} block and
its referenced \shk{instrument} block, and returns for each target a
registry row carrying the target $m/z$, its window in daltons, an
oscillation frequency, a timing-deviation interval
$[\Delta P_{\mathrm{lo}},\Delta P_{\mathrm{hi}}]$, and a minimum
acquisition time. Kind \shk{cells}; total.

\begin{remark}[Cross-block reference]
\label{rem:crossblock}
This is the only operation that reads declarative blocks other than its
own arguments: it resolves \shk{target\_list} and \shk{instrument} by name
from the AST. \Cref{rule:call} therefore passes $A$ to every denotation,
though only this operation consults it.
\end{remark}

\subsection{Namespace \shk{lavoisier.msms}}

\begin{longtable}{@{}p{0.30\textwidth}p{0.44\textwidth}p{0.18\textwidth}@{}}
\toprule
\textbf{Operation} & \textbf{Arguments} & \textbf{Kind / totality}\\
\midrule
\endhead
\shk{sebd\_search} & precursor and fragment specification & \shk{records}; partial\\
\shk{phase\_coherence} & precursor, fragment set & \shk{subharmonics}; total\\
\shk{virtual\_tensor} & ion specification, stacking dimensions & \shk{tensor}; total\\
\shk{impossible\_ions} & candidate set & \shk{impossible}; total\\
\shk{partition\_complement} & address & \shk{scalar}; total\\
\shk{transient\_contents} & transient specification & \shk{object}; total\\
\bottomrule
\end{longtable}

\subsection{Namespace \shk{lavoisier.db}}

\shk{search}, \shk{search\_massbank}, \shk{search\_gnps}, \shk{search\_mona}
query spectral reference collections
\citep{horai2010massbank,wang2016gnps} and return match records. Kind
\shk{records}; partial (network-dependent). \shk{generate} produces
records from a generative address scheme rather than a stored
collection. Kind \shk{records}; total.

\subsection{Namespace \shk{lavoisier.observe}}

\shk{partition\_field}, \shk{sentropy}, \shk{ternary\_address},
\shk{dual\_path\_validate} compute coordinate representations from
supplied observables. Kinds \shk{object} or \shk{scalar}; total.

\subsection{Namespace \shk{lavoisier.purpose}}

\shk{domain(domain, depth?)} returns the prefix set covering a named
analytical domain; \shk{match(sentropy)} returns the domains containing a
point; \shk{combine} intersects domains. Kind \shk{object}; total.

\begin{remark}[Domains are declared data, not derived]
\label{rem:domains}
The domain bounds are a fixed table shipped with the library. They are
declared constants of the implementation, and a program that relies on
them inherits whatever empirical basis they have. A program requiring
different bounds must state them explicitly rather than naming a domain.
\end{remark}

\subsection{Namespace \shk{validate}}

Validation blocks admit \shk{check\_resolution\_time(instrument,
window\_ppm)}, which reports a minimum acquisition time for the named
instrument at the given window and appends it to $L$. Validation
operations produce no bindings and cannot fail the run; they are
reporting constructs.

% =====================================================================
\section{Worked programs}
\label{sec:programs}
% =====================================================================

\subsection{Untargeted lipidomics}

\begin{lstlisting}[caption={Untargeted lipid library over five classes.}]
import lavoisier.instrument

objective PlasmaLipidomics:
    target: "predict lipidomics library for plasma LC-MS/MS"
    success_criteria: "records > 100"

instrument OrbitrapFusion:
    analyzer: "orbitrap"
    polarity: "+"
    collision_energy: 25
    mz_window: [200, 1500]

phase Design:
    classes = ["PC", "PE", "SM", "TAG", "Cer"]

phase VirtualRun:
    records = lavoisier.instrument.run_experiment(
        classes: classes,
        polarity: "+",
        analyser: "orbitrap",
        collision_energy: 25
    )
\end{lstlisting}

The \shk{Design} phase binds a list; the \shk{VirtualRun} phase references
it by bare word, resolved by \textsc{Var} (\Cref{rule:lit}). By
\Cref{thm:isolation} the binding is visible because \shk{Design} precedes
\shk{VirtualRun} in source order --- and would not be were the phases
transposed.

\subsection{Targeted acquisition with per-class ranges}

\begin{lstlisting}[caption={Per-class ranges and post-generation filters.}]
phase Design:
    panel = [
        { class: "PC",  carbons: [30, 40], db: [0, 4] },
        { class: "PE",  carbons: [32, 38], db: [0, 3] },
        { class: "TAG", carbons: [48, 56], db: [1, 6] }
    ]

phase VirtualRun:
    records = lavoisier.instrument.run_experiment(
        classes: panel,
        adducts: ["[M+H]+", "[M+Na]+", "[M+NH4]+"],
        gradient_min: 30,
        filters: {
            classes: ["PC", "PE"],
            rt: [8.0, 18.0],
            mz: [600, 900],
            db: [1, 4],
            min_intensity: 0.02
        }
    )
\end{lstlisting}

The \shk{panel} literal spans five physical lines and is reassembled into
one logical line by \Cref{def:logical}. Note that the generated scope
(three classes) and the retained scope (two classes, restricted window)
are stated separately, so the document records both what was enumerated
and what was kept (\Cref{rem:filters}).

\subsection{Timing-cell compilation}

\begin{lstlisting}[caption={Target list compiled to timing cells.}]
import lavoisier.cells

objective TargetedAcquisition:
    target: "compile plasma lipid targets to timing cells"

instrument OrbitrapXL:
    analyzer: "orbitrap"
    kappa: 1.0e12
    ref_frequency: 1.0e7
    timing_jitter: 2.0e-9

target_list PlasmaPanel:
    instrument: OrbitrapXL
    window_ppm: 5.0
    targets: [
        { name: "PC(32:0)",  mz: 734.5694 },
        { name: "PC(34:1)",  mz: 760.5851 },
        { name: "LPC(16:0)", mz: 496.3396 }
    ]

phase CompileCells:
    registry = lavoisier.cells.compile(target_list: PlasmaPanel)

validate TimingFeasibility:
    check_resolution_time(instrument: OrbitrapXL, window_ppm: 5.0)
\end{lstlisting}

Here \shk{PlasmaPanel} and \shk{OrbitrapXL} are resolved from the AST by
name (\Cref{rem:crossblock}), and \shk{instrument: OrbitrapXL} inside the
target list is a bare word (\Cref{rem:bareword}). The binding
\shk{registry} triggers clause (R2) of \Cref{def:primary}.

\subsection{Proteomics, addressing, and domain constraint}

\begin{lstlisting}[caption={Three further programs, abbreviated.}]
// (a) Tryptic peptide library
phase VirtualRun:
    records = lavoisier.instrument.run_proteomics(
        proteins: ["HSA", "HBB", "ENO1"],
        length_min: 7, length_max: 20, mc_max: 1
    )

// (b) Partition addressing
phase Compute:
    lipids = [ { name: "PC(32:0)", mass: 733.5622, adduct: "[M+H]+" } ]
    addresses = lavoisier.partition.compute_addresses(lipids)

// (c) Domain-constrained observation
phase ComputeSentropy:
    water_freqs = [1595.0, 3657.0, 3756.0]
    water_se    = lavoisier.observe.sentropy(frequencies: water_freqs)
    water_addr  = lavoisier.observe.ternary_address(sentropy: water_se,
                                                    depth: 12)

phase DomainContext:
    domain_context = lavoisier.purpose.domain(domain: "metabolomics",
                                              depth: 4)
\end{lstlisting}

In (c), \shk{water\_se} is produced in one phase and consumed in the
same phase, then \shk{domain\_context} is produced in a later phase; all
four bindings appear in the workspace in declaration order by
\Cref{thm:monotone}.

\begin{example}[Reading a program's effect set]
\label{ex:effects}
By \Cref{prop:enumerable}, the effect set of the program in
\Cref{sec:programs}.3 is
$\{\shk{lavoisier.cells.compile},\ \shk{check\_resolution\_time}\}$,
obtained by listing call heads. No execution is required, and no other
operation can be reached.
\end{example}

% =====================================================================
\section{Execution: a reference implementation and a worked experiment}
\label{sec:execution}
% =====================================================================

The preceding sections specify a language. This section reports what
happened when the specification was implemented and a program was run
against measured data. Two things are established: that the properties
proved in \Cref{sec:semantics,sec:stages} hold of a running system rather
than only of the calculus, and that a Shapeshifter program can carry an
experiment whose outcome was not known in advance --- including when that
outcome is negative.

\subsection{The reference implementation}

A reference implementation was written in Python directly from
\Crefrange{sec:lexical}{sec:acquisition}: a lexer and parser
(\Cref{sec:lexical,sec:grammar}), the kind discipline
(\Cref{sec:kinds}), the two-stage model (\Cref{sec:stages}), and eight
operations of the acquisition half (\Cref{sec:acquisition}). It is
approximately 950 lines and has no dependency beyond the standard
library. Its purpose is not performance but fidelity: every construct is
implemented as specified, so that a divergence between the specification
and the implementation is a defect in one or the other rather than a
matter of interpretation.

\begin{remark}[What a reference implementation settles]
\label{rem:reference-impl}
A specification that has never been executed can be internally consistent
and still unimplementable, because a proof obligation may be discharged in
the text by a definition that no procedure realises. Writing the
implementation from the document, without consulting any prior code, is
the cheapest available test of that. Three under-specifications were
found this way and are recorded in \Cref{rem:found-defects}.
\end{remark}

\subsection{Stage separation, observed}

\Cref{thm:stage-separation} states that compile-stage diagnostics are a
function of the source alone, and \Cref{cor:audit} that the stage can
report a program's inputs without opening them. Running the compile stage
alone on the program of \Cref{sec:experiment} produces the block
inventory, the effect set of seven operations, and the single declared
input path --- in a mean of $0.95$\,ms over twenty repetitions
($\min 0.57$, $\max 2.13$) --- while performing no file access. Execution
of the same program subsequently reads $6.5$\,MB from that path. The
separation is therefore not merely provable but observable: the ratio of
checking cost to execution cost is approximately $1:3500$, and checking
carries none of execution's effects.

\begin{remark}[The audit is the useful half]
\label{rem:audit-useful}
Of the two compile-stage outputs, the effect set
(\Cref{prop:enumerable}) and the input list
(\Cref{prop:inputs-enumerable}), the second proved the more useful in
practice. A reviewer asking ``what will this program read'' is answered
from the document; an author asking ``have I pointed this at the right
cohort'' is answered without waiting for a parse of the data.
\end{remark}

\subsection{The experiment}
\label{sec:experiment}

The program tests a claim that a coordinate transformation of the kind
defined in \Cref{def:transform-ops} is a property of the analyte rather
than of the spectrum. The data are a public reference library of
$5{,}328$ high-energy collisional dissociation MS/MS spectra covering
$727$ distinct analytes, of which $313$ were measured at all nine
normalised collision energies from $10\%$ to $80\%$. Because collision
energy is varied while analyte identity is held fixed, the library is a
designed test of exactly this question, and the answer is not known from
the data alone.

The criterion was stated in the program's \shk{objective} block before
the program was run:
\begin{equation}
\frac{\overline{d}_{\text{between}}}{\overline{d}_{\text{within}}} > 2.0
\qquad\text{and}\qquad
\lvert r_{a,\text{NCE}}\rvert < 0.3 \ \ \text{for each axis } a,
\label{eq:criterion}
\end{equation}
where $\overline{d}_{\text{within}}$ is the mean pairwise coordinate
distance among spectra of one analyte across collision energies,
$\overline{d}_{\text{between}}$ the mean distance between analyte
centroids, and $r_{a,\text{NCE}}$ the Pearson correlation of axis $a$
against collision energy. A comparison method (cosine similarity of the
raw peak lists, the established practice) and a negative control (the
same statistic under permuted analyte labels) were declared in the same
program.

\subsection{Outcome}

\begin{table}[htbp]
\centering
\small
\caption{Outcome against the criterion of \Cref{eq:criterion}, stated
before execution. Three of four conditions fail.}
\label{tab:outcome}
\begin{tabular}{@{}llrl@{}}
\toprule
\textbf{Condition} & \textbf{Required} & \textbf{Observed} & \\
\midrule
Separation ratio            & $>2.0$  & $0.930$  & fails \\
$\lvert r\rvert$, axis $S_k$ & $<0.3$  & $0.327$  & fails \\
$\lvert r\rvert$, axis $S_t$ & $<0.3$  & $0.009$  & holds \\
$\lvert r\rvert$, axis $S_e$ & $<0.3$  & $0.383$  & fails \\
\bottomrule
\end{tabular}
\end{table}

The criterion was not met (\Cref{tab:outcome}). A separation ratio below
unity states that coordinates move at least as far when collision energy
is changed at fixed analyte as when the analyte is changed: on these
data the transformation encodes spectrum state, not analyte identity.
The per-axis ratios are $1.04$, $0.49$ and $0.76$ for $S_k$, $S_t$ and
$S_e$; no axis reaches the threshold alone.

Two declared checks make the negative outcome interpretable rather than
merely disappointing. The negative control gives a ratio of $0.347$
under permuted labels, well below the $0.930$ obtained with true labels,
so the pipeline does discriminate --- there is roughly $2.7\times$ of
genuine analyte signal present, just not enough to meet the stated
threshold. The comparison method gives a mean cosine similarity of
$0.918$ between adjacent collision-energy levels, against $0.522$ across
all level pairs: the established method is stable where the criterion
required the coordinates to be stable, and the coordinates were not.

\subsection{The outcome is not a parameter choice}

\Cref{rem:totality-why} records that $\alpha$, $\beta$ and $k$ are
analyst choices which the language requires to be stated but does not
select. It is therefore worth asking whether the outcome above is an
artefact of the particular values declared in the program. Re-running the
same program over a grid of forty parameter settings --- $\alpha \in
[0,4]$, $\beta \in [0.25,4]$, $k \in [2,20]$ --- gives separation ratios
between $0.784$ and $0.968$. No setting reaches the stated threshold of
$2.0$, and none reaches even unity. Restricting the distance to each of
the seven non-empty axis subsets likewise fails, the best being $S_k$
alone at $0.962$ and the worst $S_t$ alone at $0.448$.

That the ratio is nearly flat in $\beta$ (varying by $0.002$ across a
sixteen-fold range) while falling monotonically in $\alpha$ is itself
informative: the mass term of \Cref{eq:sk} moves coordinates apart
without moving analytes apart, so weighting it more heavily degrades the
statistic. The negative outcome is a property of the transformation on
these data, not of the values chosen for the run.

\begin{remark}[What a sweep of this kind can and cannot settle]
\label{rem:sweep-limits}
Forty settings do not exhaust a three-parameter continuum, and a grid
search establishes no impossibility. What it establishes is narrower and
sufficient for the purpose: the declared setting was not unlucky, and the
gap to the threshold --- roughly a factor of two --- is not of a size that
parameter adjustment within the tested ranges closes.
\end{remark}


\begin{remark}[Why this is reported here]
\label{rem:why-report}
This paper specifies a language, not a coordinate system, and the
outcome above is evidence about the latter. It is reported because it is
evidence about the former as well: it demonstrates that a Shapeshifter
program can state a criterion, a control and a comparison in advance, and
can return a verdict against its author. A specification language whose
worked examples all succeed has not been shown to be capable of anything
in particular.
\end{remark}

\subsection{Defects found by implementing}
\label{rem:found-defects}

Three points were under-specified and are corrected in the present text.

\begin{enumerate}[label=\textbf{D\arabic*.},leftmargin=2.8em,itemsep=3pt]
\item \textbf{Analyte identity was assumed, not defined.}
      \Cref{def:scan} lists the fields a scan record carries but does not
      say which combination identifies an analyte across scans. Grouping
      on the library's name field is wrong, because that field embeds
      the retention time, which varies between acquisitions of the same
      analyte; grouping on formula alone merges isomers. The
      implementation uses formula, precursor $m/z$ and isomer index
      jointly, and the specification now requires an operation that
      groups to name its key explicitly rather than assume one.

\item \textbf{Bulk values needed a serialisation rule.} Writing a
      workspace of $5{,}298$ coordinate triples to JSON is not
      addressed by \Cref{sec:stages}. Silent truncation would corrupt any
      downstream analysis reading the artefact. The implementation
      truncates above a declared bound and records the truncation and the
      true length in the emitted object, so a consumer can detect it;
      this is the serialisation counterpart of
      \Cref{rem:null-detection}.

\item \textbf{Kind for structural literals was ambiguous at the
      boundary.} \Cref{def:classify} assigns structural kinds by
      inspection when $f=\bot$, but the implementation must decide the
      kind of a binding produced by an operation that the registry does
      not know. It raises rather than guessing, which
      \Cref{prop:failure-classification}(ii) requires but did not
      previously spell out for this case.
\end{enumerate}

\begin{remark}[Cost of the run]
\label{rem:cost}
The full program --- reading $5{,}328$ spectra comprising $133{,}220$
peaks, computing $5{,}298$ coordinate triples, and evaluating four
analyses including an $11{,}268$-pair cosine baseline --- completes in
$3.4$\,s single-threaded, of which $0.57$\,s is parsing and $2.39$\,s is
the transformation. Each workspace binding is emitted as a separate JSON
document alongside the combined record, so the artefact set is the run.
\end{remark}

% =====================================================================
\section{What the language does not provide}
\label{sec:limits}
% =====================================================================

We state the limitations precisely, because a specification language
whose boundaries are unclear cannot be relied upon.

\begin{enumerate}[label=\textbf{L\arabic*.},leftmargin=2.8em,itemsep=4pt]

\item \textbf{No argument type checking.} The kind discipline governs
      results only (\Cref{rem:no-argcheck}). An argument of wrong shape
      fails at \Cref{rule:call}, at runtime. A signature calculus for
      arguments is the natural next extension and would strengthen
      \Cref{thm:stage-separation} from structural certification to
      partial executability certification.

\item \textbf{Brackets inside string literals affect continuation.}
      \Cref{def:balance} counts brackets lexically, without string
      awareness. A string containing an unmatched bracket will extend a
      logical line incorrectly. A lexer tracking string state removes
      this.

\item \textbf{Bare words absorb misspellings.} By clause (V7) an
      unrecognised token in value position becomes a string
      (\Cref{rem:bareword}), so a misspelled field value is not an error.
      Field-level validation against a declared schema would catch it.

\item \textbf{No phase-local scope.} \Cref{thm:isolation} is a
      statement of what the language does, not an endorsement: two
      phases using the same working name silently share it. Introducing
      scope would be a compatible extension.

\item \textbf{Argument resolution is shallow.} Bare words nested inside
      arrays or objects are not resolved (\Cref{rem:resolve-depth}). This
      is deliberate but occasionally surprising.

\item \textbf{Reading is limited to \shk{mzML}.} The acquisition layer
      reads the open standard \citep{martens2011mzml,deutsch2008mzml}
      only. Vendor-native formats must be converted before a program can
      name them, and the language does not perform that conversion.

\item \textbf{No peak detection or deconvolution.} \shk{read\_mzml}
      applies intensity thresholds and a minimum-peak filter to centroided
      input; it does not perform centroiding, peak picking, isotope
      grouping, or chromatographic deconvolution. Programs consume the
      peak lists the file provides.

\item \textbf{Transformation weights are declared, not derived.} The
      parameters $\alpha$, $\beta$ and $k$ of \Crefrange{eq:sk}{eq:se},
      and the coherence threshold of \Cref{def:coherence}, are analyst
      choices. The language requires that they be stated in the source
      (\Cref{ex:threshold}); it offers no procedure for selecting them and
      no warrant that any particular value is correct.

\item \textbf{No statistical output.} There is no notion of replicate,
      no error model, and no false-discovery control. A program states an
      experiment; it does not analyse one.

\item \textbf{Domain tables are implementation data.} The
      \shk{lavoisier.purpose} domain bounds are shipped constants
      (\Cref{rem:domains}); programs depending on them inherit their
      provenance, which the language does not record.

\end{enumerate}

% =====================================================================
\section{Conclusion}
\label{sec:conclusion}
% =====================================================================

We have specified Shapeshifter: its lexical structure including the
bracket-continuation rule, its context-free grammar over six block
forms, a kind discipline in which results are tagged by their producing
operation rather than by shape, a small-step operational semantics over
configurations of environment, kind map, declaration order and log, a
two-stage compilation model separating structural checking from
execution, and an acquisition layer that brings measured data into the
same language as predicted experiments.

Four properties were proved. Declaration order determines execution and
the final configuration is unique (\Cref{thm:determinism}). Phases share
a single environment and carry no scope (\Cref{thm:isolation}). The
workspace grows monotonically, so a run is an append-only record
(\Cref{thm:monotone}). Compile-stage diagnostics are a function of the
source alone, so a program can be checked without being run and without
incurring the effects its operations would carry
(\Cref{thm:stage-separation}). To these we added three structural
observations: kind is producer-determined, so an empty result is not
retyped (\Cref{thm:kind-preservation}); the effect set of a program is
enumerable from its source, because the language admits no user-defined
abstraction (\Cref{prop:enumerable}); and a program's input files are
likewise enumerable, so the compile stage can audit what will be opened
without opening it (\Cref{prop:inputs-enumerable}, \Cref{cor:audit}).
The acquisition layer contributes one further result: measured and
predicted experiments meet in a common record kind
(\Cref{thm:record-parity}), so a document may state both and compare them
without either being converted --- while \Cref{rem:parity-limits} records
that this is a guarantee about composition, not about commensurability.

\Cref{sec:execution} reports what happened when the specification was
implemented and run. Stage separation held in practice at a checking-to-
execution cost ratio of about $1:3500$; three under-specifications were
exposed by writing the implementation and are corrected here; and the
worked experiment returned a verdict against the transformation it
tested, which is the outcome that shows the apparatus can return one.

The language is small by design, and \Cref{sec:limits} states where the
boundary lies. What it buys is a single artefact: an experiment
specification that is executable, diffable, and readable in the order it
runs. Whether that is worth the expressiveness it forgoes is a question
about practice rather than about the language, and it is settled by use
rather than by argument.

\bibliographystyle{plainnat}
\bibliography{references}

\end{document}
